% This LaTeX was auto-generated from MATLAB code.
% To make changes, update the MATLAB code and export to LaTeX again.

\documentclass{article}

\usepackage[utf8]{inputenc}
\usepackage[T1]{fontenc}
\usepackage{lmodern}
\usepackage{graphicx}
\usepackage{color}
\usepackage{hyperref}
\usepackage{amsmath}
\usepackage{amsfonts}
\usepackage{epstopdf}
\usepackage[table]{xcolor}
\usepackage{matlab}
\usepackage[paperheight=795pt,paperwidth=614pt,top=72pt,bottom=72pt,right=72pt,left=72pt,heightrounded]{geometry}

\sloppy
\epstopdfsetup{outdir=./}
\graphicspath{ {./set2sub3iun14_media/} }

\begin{document}

\begin{par}
\begin{flushleft}
\includegraphics[width=\maxwidth{63.02057200200703em}]{image_0}
\end{flushleft}
\end{par}

\matlabheading{Solutie}

\begin{matlabcode}
nodes=[-1 0 1]
\end{matlabcode}
\begin{matlaboutput}
nodes = 1x3    
    -1     0     1

\end{matlaboutput}
\begin{matlabcode}
f=@(x) cos(pi*x);
values=f(nodes)
\end{matlabcode}
\begin{matlaboutput}
values = 1x3    
    -1     1    -1

\end{matlaboutput}

\begin{par}
\begin{flushleft}
Trebuie sa gasim un spline cubic deBoor intr-o diviziune formata de n=3 puncte. Asta inseamna ca va trebui sa gasim 2 polinoame cubice ($p_i \left(x\right)=c_{i,0} +c_{i,1} \;x+c_{i,2} \;x^2 +c_{i,3} \;x^3$).
\end{flushleft}
\end{par}

\begin{par}
\begin{flushleft}
Problema spline-urilor este insuficienta ecuatiilor. Ele incearca sa egaleze nodurile de fixate cu valorile corespondente ($\left.p_i \left(x_i \right)=f_i ,\;i=0\ldotp \ldotp n-1,\mathrm{rezultand}\;\mathrm{in}\;n-1\;\mathrm{ecuatii}\right)$, dar si pentru nodurile din capatul drept (deci $p_i \left(x_{i+1} \right)=f_i ,i=0\ldotp \ldotp n-1$, rezultand tot in $n-1$ ecuatii). Pe langa acestea, ele si sa egaleze derivatele de ordinul 1 si 2 in punctele de jonctiune (${p_i }^{\prime } \left(x_{i+1} \right)={p_{i+1} }^{\prime } \left(x_{i+1} \right),$$\left.{{p_i }^{\prime } }^{\prime } \left(x_{i+1} \right)={{p_{i+1} }^{\prime } }^{\prime } \left(x_{i+1} \right),i=0\ldotp \ldotp \ldotp n-2,\mathrm{rezultand}\;\mathrm{in}\;2n-4\;\mathrm{ecuatii}\right)$
\end{flushleft}
\end{par}

\begin{par}
\begin{flushleft}
In total am obtinut $2n-2+2n-4=4n-6$ ecuatii, insa avem $4n-4$ necunoscute (4 pentru fiecare $p_{i=0\ldotp \ldotp n-1}$). De aceasta putem defini mai multe tipuri de spline sa completeze acele conditii.
\end{flushleft}
\end{par}

\begin{par}
\begin{flushleft}
Spline-ul cubic deBoor are proprietatea ca:
\end{flushleft}
\end{par}

\begin{par}
\begin{flushleft}
${{{p_1 }^{\prime } }^{\prime } }^{\prime } \left(x_2 \right)={{{p_2 }^{\prime } }^{\prime } }^{\prime } \left(x_2 \right)$ si ${{{p_{n-1} }^{\prime } }^{\prime } }^{\prime } \left(x_{n-1} \right)={{{p_{n-2} }^{\prime } }^{\prime } }^{\prime } \left(x_{n-1} \right)$
\end{flushleft}
\end{par}

\begin{par}
\begin{flushleft}
Deoarece $n=3$, problema este ca spline-ul deBoor in acest context adauga o singura informatie  $\mathrm{deoarece}\;n-1=2\;\mathrm{si}\;n-2=1$, deci a adaugat aceeasi ecuatie de 2 ori, deci vom avea un sistem subdeterminat (cu $4n-5$ ecuatii si $4n-4$ necunoscute), deci toate valorile vor depinde de o variabila. Scriem ecuatiile reiesite din spline-ul cubic deBoor:
\end{flushleft}
\end{par}

\begin{matlabcode}
syms p0 [1 4];
syms p1 [1 4];

eq01=p0(1) + p0(2)*nodes(1) + p0(3)*nodes(1)^2 + p0(4)*nodes(1)^3 == values(1), ...
    eq02=p1(1) + p1(2)*nodes(2) + p1(3)*nodes(2)^2 + p1(4)*nodes(2)^3 == values(2)
\end{matlabcode}
\begin{matlabsymbolicoutput}
eq01 = 

\hskip1em $\displaystyle p_{01} -p_{02} +p_{03} -p_{04} =-1$
eq02 = 

\hskip1em $\displaystyle p_{11} =1$
\end{matlabsymbolicoutput}
\begin{matlabcode}
eq11=p0(1) + p0(2)*nodes(2) + p0(3)*nodes(2)^2 + p0(4)*nodes(2)^3 == values(2), ...
    eq12=p1(1) + p1(2)*nodes(3) + p1(3)*nodes(3)^2 + p1(4)*nodes(3)^3 == values(3)
\end{matlabcode}
\begin{matlabsymbolicoutput}
eq11 = 

\hskip1em $\displaystyle p_{01} =1$
eq12 = 

\hskip1em $\displaystyle p_{11} +p_{12} +p_{13} +p_{14} =-1$
\end{matlabsymbolicoutput}
\begin{matlabcode}
eq2=p0(2) + 2*p0(3)*nodes(2) + 3*p0(4)*nodes(2)^2 == p1(2) + 2*p1(3)*nodes(2) + 3*p1(4)*nodes(2)^2
\end{matlabcode}
\begin{matlabsymbolicoutput}
eq2 = 

\hskip1em $\displaystyle p_{02} =p_{12} $
\end{matlabsymbolicoutput}
\begin{matlabcode}
eq3=2*p0(3) + 6*p0(4)*nodes(2) == 2*p1(3) + 6*p1(4)*nodes(2)
\end{matlabcode}
\begin{matlabsymbolicoutput}
eq3 = 

\hskip1em $\displaystyle 2\,p_{03} =2\,p_{13} $
\end{matlabsymbolicoutput}
\begin{matlabcode}
eq4=6*p0(4)==6*p1(4) % ecuatia adaugata de spline-ul deBoor
\end{matlabcode}
\begin{matlabsymbolicoutput}
eq4 = 

\hskip1em $\displaystyle 6\,p_{04} =6\,p_{14} $
\end{matlabsymbolicoutput}
\begin{matlabcode}
[p01,p02,p03,p04,p11,p12,p13]=solve([eq01 eq02 eq11 eq12 eq2 eq3 eq4], [p0(1),p0(2),p0(3),p0(4),p1(1),p1(2),p1(3)])
\end{matlabcode}
\begin{matlabsymbolicoutput}
p01 = 

\hskip1em $\displaystyle 1$
p02 = 

\hskip1em $\displaystyle -p_{14} $
p03 = 

\hskip1em $\displaystyle -2$
p04 = 

\hskip1em $\displaystyle p_{14} $
p11 = 

\hskip1em $\displaystyle 1$
p12 = 

\hskip1em $\displaystyle -p_{14} $
p13 = 

\hskip1em $\displaystyle -2$
\end{matlabsymbolicoutput}
\begin{matlabcode}
syms alfa;
c=[np01 subs(np02, p1(4), alfa) np03 subs(np04, p1(4), alfa); np11 subs(np12, p1(4), alfa) np13 alfa]
\end{matlabcode}
\begin{matlabsymbolicoutput}
c = 

\hskip1em $\displaystyle \left(\begin{array}{cccc}
1 & -\textrm{alfa} & -2 & \textrm{alfa}\\
1 & -\textrm{alfa} & -2 & \textrm{alfa}
\end{array}\right)$
\end{matlabsymbolicoutput}

\begin{par}
\begin{flushleft}
Observam ca variabilele rezultat depind de $p_{14}$.
\end{flushleft}
\end{par}

\begin{matlabcode}
function S=evalsplinecubic(x,nodes,values,c)
    S=zeros(1, length(x));
    for i=1:length(x)
         if x(i)==nodes(end)
             S(i)=values(end);
         else
             poz=find(nodes>x(i),1)-1;
             X=x(i);
             S(i)=c(poz,1)+X*(c(poz,2)+X*(c(poz,3)+X*c(poz,4)));
         end
    end
end
\end{matlabcode}


\vspace{1em}
\begin{matlabcode}
x=linspace(-1, 1, 1000);
clf; hold on; grid on;
plot(nodes, values, 'ok', 'markerfacecolor', 'r', 'markersize', 8)
plot(x, f(x))
for i=-1:0.2:1
    cn=double(subs(c,alfa,i));
    S=evalsplinecubic(x,nodes,values,cn);
    plot(x, S);
end
legend("nodes", "f")
\end{matlabcode}
\begin{center}
\includegraphics[width=\maxwidth{56.196688409433015em}]{figure_0.eps}
\end{center}

\begin{par}
\begin{flushleft}
Deci se observa ca nu obtinem doar o solutie, ci o infinitate, deci din acest motiv, nu este un spline cubic propriu zis.
\end{flushleft}
\end{par}

\end{document}
